\section{Conclusiones}

El presente trabajo de grado tuvo como objetivo general desarrollar un prototipo web que apoyara la comprensión de los contenidos de la asignatura Fundamentos de Interpretación y Compilación de Lenguajes de Programación (FLP) de la Universidad del Valle, sede Tuluá, mediante la visualización de estructuras internas y el uso de lenguajes orientados por sintaxis. A partir del diagnóstico de dificultades conceptuales, el diseño y la implementación del prototipo FLP Viewer, y su posterior evaluación con estudiantes de la asignatura, se considera que este objetivo fue alcanzado.

\subsection{Cumplimiento de los objetivos específicos}

El primer objetivo específico, orientado a identificar las dificultades conceptuales que enfrentan los estudiantes de la asignatura, se desarrolló en el Capítulo~\ref{cap3} mediante la aplicación de encuestas a 36 estudiantes y a los docentes responsables de los distintos niveles de la asignatura. El diagnóstico permitió confirmar que las dificultades no se concentran en la comprensión teórica general, sino específicamente en el segundo resultado de aprendizaje del microcurrículo, relacionado con ambientes de ejecución, estado y ligadura de variables, donde el 77.8\,\% de los estudiantes reportó dificultades explícitas. Este hallazgo orientó directamente la selección de los cinco indicadores de logro que sirvieron como base conceptual para el prototipo.

El segundo objetivo, referido a la caracterización de los fundamentos teóricos y técnicos sobre gramáticas formales, análisis sintáctico y árboles de sintaxis abstracta, se desarrolló en el Capítulo~\ref{cap2} mediante la construcción de un marco teórico que integró los fundamentos técnicos de la interpretación de lenguajes, las teorías del aprendizaje aplicadas a la enseñanza de compiladores y los principios de usabilidad en herramientas educativas interactivas. La búsqueda sistemática de literatura documentada en el Anexo~\ref{annex:slr} permitió, adicionalmente, identificar los vacíos que fundamentaron el aporte específico de este trabajo frente a las herramientas existentes.

El tercer objetivo, relacionado con el diseño de la arquitectura web y las interfaces de usuario, se desarrolló en el Capítulo~\ref{cap4}, donde se especificaron los requisitos funcionales y no funcionales del sistema y se definieron las decisiones de arquitectura, las tecnologías empleadas y las estrategias de visualización orientadas a mantener en un único entorno las representaciones que el estudiante debe relacionar para comprender el proceso de interpretación.

El cuarto objetivo, correspondiente a la implementación de un prototipo funcional, se desarrolló en el Capítulo~\ref{cap5}. El sistema implementado permite al estudiante definir gramáticas en notación BNF, generar y visualizar de forma interactiva el AST resultante, y observar los cambios en el ambiente de ejecución durante la evaluación paso a paso de un programa, cumpliendo con los ocho requisitos funcionales especificados en el Anexo~\ref{annex:rf}.

El quinto objetivo, relacionado con la evaluación de la usabilidad y funcionalidad del prototipo mediante pruebas con estudiantes y el docente responsable, se desarrolló en el Capítulo~\ref{cap6}.

\subsection{Aportes del trabajo}

Frente a los vacíos identificados en la revisión de antecedentes (Capítulo~\ref{cap2}), este trabajo aporta un prototipo que permite al estudiante definir su propia gramática en notación BNF en lugar de operar sobre un lenguaje predefinido, e integra en un mismo entorno la edición de la gramática, la visualización del AST y la inspección del ambiente de ejecución, un flujo que no se encontró combinado en ninguna de las herramientas revisadas. Adicionalmente, el trabajo aporta evidencia empírica sobre el uso de una herramienta de este tipo en un curso de lenguajes de programación de un programa de ingeniería en una universidad latinoamericana, un contexto poco representado en la literatura identificada.

\subsection{Limitaciones}

La evaluación del prototipo presenta limitaciones que deben considerarse al interpretar los resultados presentados en el Capítulo~\ref{cap6}. El tamaño de la muestra (12 estudiantes de una única cohorte) no permite generalizar los hallazgos a otras cohortes o instituciones. La evaluación se restringió, además, a la usabilidad y a la utilidad pedagógica percibida, dejando por fuera mediciones del impacto del prototipo en el desempeño académico o en la retención del aprendizaje a largo plazo, tal como se estableció en el alcance del trabajo. Finalmente, el prototipo continúa dependiendo de Racket y EOPL como sustrato de ejecución, por lo que la carga cognitiva asociada a la sintaxis de Racket señalada en el Capítulo~\ref{cap3} no desaparece, sino que se busca mitigar mediante la representación visual complementaria.
